\section{官方幻灯片证据链：从可验证数据到稳定探索}

本节按官方 deck 的顺序重建整套后训练逻辑。课程不是把“Agent”定义为多调用几次工具，而是把训练对象从单轮回答扩展为环境中的行动轨迹：模型读取状态、选择工具、改变环境，再由 verifier 判断是否达到可验证目标。读图时要持续区分人类偏好、程序化正确性、探索多样性和计算成本，因为四者并不会自动同步改善。

\subsection{Agentic shift：从 human preference 对齐到 environment feedback 对齐}

slides 1--6 先比较聊天模型与 Agentic 模型。早期 chat post-training 用 SFT、PPO、GRPO 等方法吸收偏好数据，目标是生成更令人满意的回答；Agent 模型还要与 sandbox、工具和环境状态交互，以 rubric、ground truth、unit test 等信号最大化正确性。讲者随后把课程拆成三步：获得可验证数据、定义“智能”的评估、找到把数据有效喂给模型的 recipe。

\lecturefigure{slides-images/slide-001.png}{官方 slide 1：Post-Training Verifiable Agents}
\lecturefigure{slides-images/slide-002.png}{官方 slide 2：human preference aligned 与 environment feedback aligned}
\lecturefigure{slides-images/slide-003.png}{官方 slide 3：早期 chat model 的偏好后训练}
\lecturefigure{slides-images/slide-004.png}{官方 slide 4：可验证任务、环境交互与聊天任务的差异}
\lecturefigure{slides-images/slide-005.png}{官方 slide 5：sandbox、tools 与 verifier 驱动的 Agentic model}
\lecturefigure{slides-images/slide-006.png}{官方 slide 6：数据、智能定义与训练 recipe 三步框架}

\begin{knowledgebox}{读图：可验证奖励不是人类偏好的替代品}
slide 2 写的是 ``in addition to human preference''。代码通过测试不代表修改可维护，数学答案正确不代表解释适合用户，网页任务完成也不代表没有越权。生产模型通常需要联合目标：任务正确、交互可用、行为安全、成本可控。可验证信号的价值在于把某些结果从主观评价转成可重复执行的检查，而不是把所有价值判断压缩成一个测试。
\end{knowledgebox}

\textbf{课堂提示：}讲者将 environment feedback alignment 作为 Agentic 模型的分界线。环境不是一段背景文字，而是会随动作变化的状态；工具也不是固定格式，而是读写环境的接口。若训练数据只保存最终答案，不保存动作与反馈，模型无法学习哪一步导致成功或失败。

\subsection{可验证训练数据：Environment、Tools、Verifier 三元组}

slides 8--11 定义数据结构。Environment 可以是代码仓库、浏览器、数据库或销售流程；Tools 是查询或修改状态的 API；Verifier 可以是单元测试、数学检查器、DOM 脚本或 ground-truth state。一个好数据集必须覆盖多种环境、工具和 verifier，并控制 verifier 的 false positive/false negative。困难任务尤其容易出现“答案其实正确但未被接受”或“错误轨迹碰巧过测试”的问题。

\lecturefigure{slides-images/slide-008.png}{官方 slide 8：Environment、Tools 与 Verifier 三个训练对象}
\lecturefigure{slides-images/slide-009.png}{官方 slide 9：三轴多样性决定训练数据覆盖}
\lecturefigure{slides-images/slide-010.png}{官方 slide 10：API、数据库、浏览、代码工具与多类环境/verifier}
\lecturefigure{slides-images/slide-011.png}{官方 slide 11：verifier false positive 与 false negative}

\begin{knowledgebox}{术语消化：状态、动作与 verifier}
状态 $s_t$ 表示第 $t$ 步环境可观测信息，动作 $a_t$ 是模型生成的工具调用或文本操作，环境返回新状态 $s_{t+1}$。轨迹写作 $\tau=(s_0,a_0,s_1,\ldots,s_T)$。verifier $V(\tau)$ 根据终局状态或整个过程返回奖励。若 $V$ 只检查最终文本，模型可能通过破坏环境、修改测试或走危险捷径得分；因此高风险任务还要验证权限、过程不变量和副作用。
\end{knowledgebox}

\begin{warningbox}{Verifier quality 决定奖励是否可信}
同一个数学值可能写成 $1/2$ 或 $2/4$；若题目要求最简形式，二者待遇不同。verifier 必须把题意编码准确。false negative 会惩罚正确的多样解法，压缩探索；false positive 会奖励错误策略，引发 reward hacking。\textbf{老师强调：}规模不能弥补 verifier 定义错误，困难 prompt 上反而更容易放大错误。
\end{warningbox}

\subsection{Holistic evaluation：benchmark、harness 与智能定义}

slides 13--18 把评估从单一 benchmark 扩展为五个维度：工具、use case、模糊指令、软件系统和鲁棒性。off-the-shelf benchmark 测任务，agent harness 决定模型怎样探索 sandbox，focused unit eval 检查结构输出、条件指令和长轨迹。课程进一步提出 benchmark quality 的三个轴：hardness、diversity 和 separability。选择评估套件，就是在操作性地定义“智能”。

\lecturefigure{slides-images/slide-013.png}{官方 slide 13：Agent 评估需要量化追踪的五类期望}
\lecturefigure{slides-images/slide-014.png}{官方 slide 14：benchmark、agent harness 与 focused unit evaluation}
\lecturefigure{slides-images/slide-015.png}{官方 slide 15：任务与多种 harness 组成 holistic measurement}
\lecturefigure{slides-images/slide-016.png}{官方 slide 16：多任务、多 harness、多工具共同定义智能}
\lecturefigure{slides-images/slide-017.png}{官方 slide 17：任务集合是否足以成为智能度量}
\lecturefigure{slides-images/slide-018.png}{官方 slide 18：hardness、diversity 与 separability}

\begin{knowledgebox}{读图：Harness swap 在测什么}
同一模型接 OpenHands、bash-only MiniSWE 或不同搜索 harness，结果会显著变化。差异可能来自模型，也可能来自 scaffold 提供的工具、状态压缩、错误恢复和提示。固定 harness 只能测“模型+系统”的组合；交换 harness 才能估计模型对接口变化的鲁棒性。评估报告应同时记录 model、harness、tool version 和环境镜像，避免把系统升级误当成模型能力提升。
\end{knowledgebox}

\begin{importantbox}{Benchmark quality 的三个轴}
Hardness 要让当前模型既非全对也非全错，才能提供学习和比较信号；Diversity 要覆盖不同环境、工具和解法；Separability 要能稳定区分能力不同的模型。难但无稳定 verifier 的题不一定是好 benchmark，题量大但同质也不代表 holistic。Agent 评估还应报告成本、轨迹长度和失败类型。
\end{importantbox}

\subsection{先模仿再探索：Light SFT 与 RL 的职责边界}

slides 20--23 把训练解释为在环境中采样多条 attempt，并强化正确且有价值的轨迹。SFT 先模仿专家示范，目的是降低采到荒谬动作的概率，使后续 rollout 在有限 compute 下可用；RL 再从当前策略探索不同轨迹。示范应多样，SFT 应轻，因为过度拟合单一示范会让采样分布变窄，RL 无法看到足够多的可替代策略。

\lecturefigure{slides-images/slide-020.png}{官方 slide 20：多 attempt、工具、verifier 与采样成本}
\lecturefigure{slides-images/slide-021.png}{官方 slide 21：SFT imitate 与 RL explore}
\lecturefigure{slides-images/slide-022.png}{官方 slide 22：light/diverse SFT 为 RL 提供有意义的 rollout}
\lecturefigure{slides-images/slide-023.png}{官方 slide 23：SFT 排除低质量，RL 强化智能}

\begin{knowledgebox}{读图：SFT 与 RL 不是“基础版/高级版”}
SFT 优化示范轨迹的似然，适合注入格式、工具协议和初始策略；RL 需要在线或近在线采样，再根据 reward/advantage 更新。若零样本策略几乎不会成功，RL 预算会浪费在无意义轨迹上；若 SFT 太强且示范单一，策略熵降低，RL 又缺少探索空间。二者的接口是 rollout 分布，而不是固定训练轮数。
\end{knowledgebox}

\textbf{老师强调：}“SFT is just there”是强调其在本讲 recipe 中的功能边界，不是贬低 SFT。高质量示范仍决定工具语义和安全起点；但真正超越示范、发现新轨迹，需要当前策略在环境中探索并接受 verifier 反馈。

\subsection{Train longer：稳定性、策略偏差与 entropy balance}

slides 24--30 给出 good RL 的第一支柱。训练更久不是增加 step 数，而是控制 reward 上升与 entropy 下降的交换。off-policy 数据会使更新偏离当前策略，产生 biased update；clip 过于对称或过紧，会阻止低概率但有价值 token 增长；显式 entropy loss 可鼓励分布保持多样。课程引用 on-policy baseline、DAPO、Skywork 等结果说明这些干预如何影响 generalization。

\lecturefigure{slides-images/slide-024.png}{官方 slide 24：train longer、train harder、sample better 三支柱}
\lecturefigure{slides-images/slide-025.png}{官方 slide 25：entropy-reward tradeoff 与训练时长}
\lecturefigure{slides-images/slide-026.png}{官方 slide 26：减少 biased update 与平衡 clipping}
\lecturefigure{slides-images/slide-027.png}{官方 slide 27：on-policy 更新、entropy control 与 generalization}
\lecturefigure{slides-images/slide-028.png}{官方 slide 28：非对称 clipping 为低概率 token 留出增长空间}
\lecturefigure{slides-images/slide-029.png}{官方 slide 29：显式 entropy regularization 路径}
\lecturefigure{slides-images/slide-030.png}{官方 slide 30：entropy loss 的训练曲线}

\begin{knowledgebox}{术语消化：on-policy、off-policy、clipping 与 entropy}
On-policy 使用当前策略采样的数据更新当前策略；off-policy 使用旧策略或其他策略数据，复用效率更高但 importance ratio 可能带来偏差。Clipping 限制新旧策略概率比，防止单次更新过强；上下界可非对称，让原本低概率的好动作有更大增长空间。Entropy 衡量动作分布不确定性，下降过快意味着策略过早集中，可能丢失可泛化探索。
\end{knowledgebox}

\begin{warningbox}{Entropy 高低都不是独立目标}
高 entropy 可能只是随机和冗长，低 entropy 也可能是模型已学会确定策略。课程关注的是与 reward、任务难度和 generalization 联合观察。显式 entropy bonus 若过大，会牺牲正确性；若完全不控制，训练可能快速塌缩。图中曲线支持特定实验干预，不能直接给出所有 Agent 任务的通用系数。
\end{warningbox}

\subsection{Train harder：任务应处于当前策略的可学习区间}

slides 31--34 解释第二支柱。不能把最难题一股脑加入训练；若当前策略几乎从不成功，reward 全为零，梯度没有方向。也不能只训练已饱和的简单题。合适难度要求策略置信与 advantage/reward 之间保留有意义相关。课程用数学 benchmark 曲线展示过易、合适和过难数据，并讨论提高困难 prompt 权重来改善 learning signal。

\lecturefigure{slides-images/slide-031.png}{官方 slide 31：难度、策略概率与 advantage 的相关性}
\lecturefigure{slides-images/slide-032.png}{官方 slide 32：简单、适中与超难数学数据的训练差异}
\lecturefigure{slides-images/slide-033.png}{官方 slide 33：改善困难任务 learning signal 的干预}
\lecturefigure{slides-images/slide-034.png}{官方 slide 34：对更难 prompt 提高奖励权重}

\begin{knowledgebox}{读图：难度课程应随模型移动}
同一道题对不同 checkpoint 难度不同。可用 pass@K、成功率或 verifier reward 分布估计当前可学习区间：成功率接近 100\% 的题贡献有限，接近 0 的题又缺少正例；中间区域通常最有信息。随着策略提升，curriculum 要动态加入更难题，并保留部分旧题监控遗忘。难度加权也要防止模型只优化高权重窄领域。
\end{knowledgebox}

\textbf{课堂提示：}讲者把“hard but not impossible”与稳定训练绑定，而不是只追求 benchmark headline。Agent 环境中还要考虑失败成本：某些探索虽有学习价值，却可能产生危险副作用，因此困难任务应在 sandbox 中设计，并由权限和过程 verifier 约束。

\subsection{Sample better：扩大计算必须配合选择与置信度}

slides 35--38 给出第三支柱。增加每题采样数可获得更多推理路径，但只有高质量反馈和选择器才能把 compute 转成训练信号。GenSelect 把多个候选交给选择模型挑最佳；DeepConf 只在置信度高于阈值的候选上做 majority voting；beam-style reasoning 则在中间步骤搜索。最终 recipe 把 expert/synthetic/self-distilled SFT 与 train longer/harder/sample better 的 RL 组合起来。

\lecturefigure{slides-images/slide-035.png}{官方 slide 35：parallel reasoning、beam search 与 sample better}
\lecturefigure{slides-images/slide-036.png}{官方 slide 36：GenSelect 从多个采样答案中选择最佳}
\lecturefigure{slides-images/slide-037.png}{官方 slide 37：DeepConf 用置信阈值筛选后投票}
\lecturefigure{slides-images/slide-038.png}{官方 slide 38：SFT 数据与三类 RL recipe 总结}

\begin{knowledgebox}{读图：Best-of-N 的收益来自覆盖与选择器}
若单样本成功率为 $p$，独立采样 $N$ 次至少一次成功的概率为 $1-(1-p)^N$；但系统不知道哪条成功，仍需 verifier 或 selector。候选高度相关时，独立假设会高估收益；selector 偏差也会选中更自信、更长却错误的答案。报告 sampling 方法时应同时给准确率、token/latency 成本、候选多样性和选择器错误。
\end{knowledgebox}

\begin{warningbox}{更多样本不等于更多有效探索}
重复同一路径只增加账单。temperature、prompt、tool policy、search branching 和模型 ensemble 都会改变多样性。DeepConf 的阈值过滤可能提高投票质量，也可能丢掉罕见正确路径；GenSelect 的 selector 需要独立校准。课程的结论是“sample better”，不是“无限 sample”。
\end{warningbox}

\subsection{开放式 collection：把环境、评估与算法放进同一公共闭环}

slides 40--44 从单个 recipe 转向生态系统。第一步众包高质量环境、verifier 和评估；第二步研究 benchmark 的强弱并形成更 holistic 的智能定义；第三步发布能稳定训练、更难课程和多样响应的算法。最终开放问题借鉴人类学习：跨环境学习、重点课程、自主探索与教师示范如何平衡，以及不同智能定义如何比较。

\lecturefigure{slides-images/slide-040.png}{官方 slide 40：collection、智能定义与训练算法三步议程}
\lecturefigure{slides-images/slide-041.png}{官方 slide 41：众包 environments、evals 与 recipes}
\lecturefigure{slides-images/slide-042.png}{官方 slide 42：分析 benchmark 并用于 holistic RL scale-up}
\lecturefigure{slides-images/slide-043.png}{官方 slide 43：稳定、困难、多样的算法进入共享 collection}
\lecturefigure{slides-images/slide-044.png}{官方 slide 44：环境、课程、探索/教师与智能定义开放问题}

\begin{importantbox}{共享 collection 的最小治理要求}
每个环境需要版本、许可证、工具权限、可复现镜像和 verifier 测试；每个 benchmark 需要污染、难度、分离度和 harness 记录；每个 recipe 需要模型、数据、采样预算、随机种子与失败日志。没有这些元数据，collection 会变成不可比较的任务堆，而不是可累积的科学基础设施。
\end{importantbox}

\textbf{讲义提醒：}用人类学习作类比能提出好问题，却不能直接推出算法。人类环境、记忆、社会反馈和身体约束远比当前 sandbox 复杂。课程将它作为研究议程：建立可实验的环境集合，明确测量，再逐步验证 curriculum、exploration 和 imitation 的组合，而不是凭类比宣布答案。

